\begin{figure}[ht]
\centering
\ThesisBox{0.24\textwidth}{\textbf{Cas retenu}\\[0.4em]\texttt{6O3MzPeomqs-H4-0h\_iDWME}\\[0.4em]Source: 235.0~s\\GT: 01:41--02:48\\Pr\'edit: 01:41--02:33\\Txt IoU: 0.768\\Temp IoU: 0.783\\Score: 0.776}
\hfill
\ThesisBox{0.33\textwidth}{\textbf{Detailed Visual Timeline}\\[0.3em]No detailed visual notes selected.}
\hfill
\ThesisBox{0.33\textwidth}{\textbf{Scenes + Dialogue}\\[0.3em]No scene summary available.\\[0.5em]\textbf{Extrait choisi}\\[0.3em]So the biggest reliable difference that's been\\ documented between males and females, and this\\ becomes even larger in egalitarian societies, by\\ the way, is orientation of interest. So women are\\ higher in negative emotion and they're more\\ agreeable. And agreeableness is both compassion\\ and politeness.}

\vspace{0.8em}
\ThesisBox{0.94\textwidth}{%
\textbf{Lecture du cas}\\[0.4em]
\TimelineBar{R\'ef\'erence GT}{0.433}{0.282}{MemoBlue}\\[0.6em]
\TimelineBar{Span pr\'edit}{0.431}{0.224}{MemoGreen}\\[0.6em]
Le mod\`ele retrouve ici le bon noyau argumentatif et reste proche des bornes de r\'ef\'erence, ce qui explique le double recouvrement textuel et temporel \'elev\'e.
}
\caption{Exemple concret de transformation d'une vid\'eo en contexte structur\'e, puis en span s\'electionn\'e par le mod\`ele.}
\label{fig:worked-example}
\end{figure}
